\documentclass[11pt]{article}

\usepackage[margin=1in]{geometry}
\usepackage{amsmath,amssymb}
\usepackage{booktabs}
\usepackage{enumitem}
\usepackage[hidelinks]{hyperref}
\usepackage{xcolor}

\newcommand{\method}{HLC-SG}
\newcommand{\todo}[1]{\textcolor{red}{[TODO: #1]}}
\newcommand{\rzero}{$r@0$}
\newcommand{\condres}{conditional resurrection}
\newcommand{\conduc}{conditional AUC}

\title{Durable Unlearning Needs Immediate-Matched Survival Evaluation}
\author{Harshad317 \and Codex Research Draft}
\date{\today}

\begin{document}
\maketitle

\begin{abstract}
Unlearning methods for language models are often evaluated immediately after intervention, but deployment exposes models to subsequent benign or answer-bearing updates that can resurrect forgotten information.
We study this gap through a local-first survival-curve harness for durable unlearning and a half-life-control objective, \method, that penalizes post-relearning resurrection.
The central experimental point is methodological: survival AUC can overstate durability when a method simply forgets more at time zero.
We therefore introduce immediate-matched comparisons and conditional resurrection metrics that measure only items forgotten at $t=0$.
In current local experiments, this protocol converts an initially promising benign HLC signal into a stricter diagnostic result: HLC variants beat weak baselines but do not yet establish a durable-unlearning win over the strongest immediate-matched baselines.
Seeded Qwen3 long-horizon results are pending and will determine whether the tuned HLC configuration supports a positive method claim or remains a negative diagnostic case study.
\end{abstract}

\section{Introduction}
Language-model unlearning asks for a model that no longer exposes a targeted set of facts while preserving general utility.
Most evaluations ask whether the model forgets immediately after an unlearning update.
That is necessary but incomplete.
A deployed model may later be fine-tuned, adapted, or otherwise updated on data that is benign, adjacent to the retained distribution, or even answer-bearing.
If forgotten facts quickly reappear under such updates, the original intervention was not durable.

This paper starts from a simple failure mode: a method can look durable because it forgets more at $t=0$, not because the forgotten items are harder to relearn.
Survival curves and survival AUC are then confounded by immediate forget strength.
We address this by measuring durability conditionally: among items that are forgotten at $t=0$, what fraction resurrect after a fixed number of relearning steps?

We make three contributions.
First, we define a reproducible survival-curve harness for durable unlearning with fixed thresholds, relearn pools, seeded evaluation, retain utility, refusal checks, and long-horizon artifacts.
Second, we introduce immediate-matched durability evaluation: baselines are selected or swept until their immediate resurrection rate, utility, and refusal behavior are comparable to the method under test.
Third, we implement and diagnose Relearning Half-Life Control with surrogate gradients, a method that explicitly penalizes resurrection after simulated relearning.

The current draft is intentionally conservative.
Early benign survival curves made \method{} look strongest, but answer-bearing stress and immediate-matched comparisons rejected a broad durability claim for the implemented HLC variants.
The active Qwen3-0.6B seeded long-horizon run is the next decision gate for a tuned HLC configuration.
Until that table is locked, the main stable claim is about evaluation rigor, not state-of-the-art performance.


\section{Problem Setup}
Let $F$ be the forget set and $R$ the retain set.
An unlearning method starts from a model $\theta_{\mathrm{full}}$ trained on both sets and produces an unlearned model $\theta_0$.
The goal is not only that $\theta_0$ fails to recover answers from $F$ immediately, but that facts in $F$ remain suppressed after subsequent relearning updates.

For each forget item $i \in F$, the evaluator uses prompt variants $p_{i,j}$, a target answer $a_i$, optional aliases, and a neutral answer $n_i$.
The primary score is a target log-likelihood margin,
\[
  m_i(\theta) = \max_j \left[\log p_{\theta}(a_i \mid p_{i,j}) - \log p_{\theta}(n_i \mid p_{i,j})\right].
\]
Thresholds $\tau_i$ are calibrated from a retain-only oracle using fixed quantiles.
An item is resurrected under the margin metric when
\[
  m_i(\theta) > \tau_i.
\]

After unlearning, the evaluator applies relearning updates on a pool $B$ for horizons
\[
  t \in \{0,1,2,4,8,16,32,64,128,256,512\}.
\]
The model after $t$ relearning steps is $\theta_t$.
The basic resurrection fraction is
\[
  r_t = \frac{1}{|F|}\sum_{i \in F} \mathbf{1}[m_i(\theta_t) > \tau_i],
\]
and survival is $S_t = 1-r_t$.

This setup separates three different requirements:

\begin{enumerate}[leftmargin=*]
  \item immediate forgetting: low $r_0$;
  \item retained utility: low retain negative log-likelihood and low refusal;
  \item durability: low resurrection growth under relearning.
\end{enumerate}


\section{Relearning Half-Life Control}
\method{} is an HLC-style unlearning objective that trains the unlearned model to remain forgotten after simulated inner relearning updates.
The implementation uses LoRA adapters for local GPU feasibility and simulates $K$ inner SGD steps on a relearn pool.
At selected inner times, it penalizes resurrection of forget-set items.

At a high level, each outer step uses:

\begin{enumerate}[leftmargin=*]
  \item a forget batch for negative preference or immediate forget pressure;
  \item a retain batch for utility preservation;
  \item an inner sequence of relearning batches;
  \item resurrection penalties after one or more simulated relearning steps;
  \item a top-$k$ utility KL term against the full model.
\end{enumerate}

The current strongest Qwen tuning candidate uses:
\[
K=4,\quad \eta_{\mathrm{inner}}=2\cdot 10^{-5},\quad
\eta_{\mathrm{outer}}=2\cdot 10^{-5},
\]
\[
\lambda_0=1,\quad \lambda_K=4,\quad \lambda_R=2,\quad
\lambda_{\mathrm{growth}}=1,
\]
with resurrection penalties at inner steps $1,2,4$ and margin-growth target $1$.

The method is best understood as a local diagnostic objective rather than a completed SOTA algorithm.
Its purpose is to test whether explicitly controlling post-relearning resurrection can move the conditional survival curve after immediate matching.
The current evidence shows that this objective improves some short-horizon signals, but the final Qwen seeded long-horizon comparison is still needed before making a positive claim.


\section{Immediate-Matched Survival Evaluation}
The evaluation protocol is designed to prevent a common overclaim: confusing stronger immediate forgetting with durable forgetting.

\paragraph{Immediate matching.}
For each method under test, candidate baselines are swept and selected so that their $r@0$ is within a tolerance band of the reference method.
The local default tolerance is $0.05$.
Candidates are rejected when retain NLL or refusal rate is outside the allowed band.
This creates a fairer comparison: both methods start with approximately the same number of resurrected forget items.

\paragraph{Conditional resurrection.}
Let
\[
  F_0 = \{i \in F : m_i(\theta_0) \le \tau_i\}
\]
be the set of items forgotten at $t=0$.
Conditional resurrection at horizon $T$ is
\[
  c_T = \frac{1}{|F_0|}\sum_{i \in F_0}\mathbf{1}[m_i(\theta_T) > \tau_i].
\]
This directly measures whether forgotten items stay forgotten.

\paragraph{Conditional survival AUC.}
Conditional survival is $1-c_t$.
We compute running AUC over the fixed horizon grid.
The most informative horizons in the current harness are $t=32$, $t=128$, and $t=512$.

\paragraph{Stress pools.}
Benign relearn pools are useful for smoke tests, but answer-bearing pools are decisive for durability.
The harness includes declarative, paraphrased QA, and mixed answer-bearing pools, plus held-out versions with different wording.
Held-out pools prevent tuning only to one narrow relearning phrasing.

\paragraph{Pass condition.}
A positive HLC claim requires lower conditional resurrection or higher conditional survival AUC than the strongest immediate-matched baseline, while retaining similar utility and low refusal.
If HLC only wins on $r@0$, the result is not a durability win.


\section{Experimental Protocol}
The implementation is local-first but preserves GPU-ready command contracts.
The early MVP used tiny GPT-2 and GPT-2 LoRA smoke runs to validate the harness.
The current active GPU lane uses Qwen3-0.6B with LoRA adapters on two V100 16GB GPUs.

\paragraph{Datasets.}
The local harness uses a TOFU-style synthetic biography fixture with a forget split and retain split.
Each forget item has prompt variants, a target answer, aliases where available, and a neutral answer.

\paragraph{Threshold calibration.}
Thresholds are calibrated from a retain-only oracle checkpoint.
The evaluator stores quantiles $q90$, $q95$, and $q99$ with global threshold $\tau_{\mathrm{global}}=0$.
The default reported metric uses $q95$.

\paragraph{Baselines.}
Implemented baselines include gradient ascent with retain CE, NPO, strong NPO, NPO with paraphrases, syntactic diversification, and immediate-matched constrained gradient ascent.
The decisive Qwen comparison currently uses constrained gradient ascent because Strong NPO did not reach HLC's immediate forget level in the Qwen sweep.

\paragraph{HLC variants.}
The draft covers K1/K2 GPT-2 smoke variants, prior K2/K4 local stress diagnostics, and the active Qwen K4 stress-tuned candidate.
The final Qwen K4 long-horizon table is pending seeded rerun completion.

\paragraph{Reproducibility.}
The survival evaluator writes CSV and JSONL artifacts for each horizon, including survival curves, retain utility, resurrection examples, and generation logs.
All result tables in this draft should be traceable to `runs/reports/` artifacts.


\section{Results}
This section is split into locked local diagnostics and pending Qwen results.

\subsection{Benign Curves Overstate HLC Durability}

On benign relearn pools, HLC-SG appeared strongest in the local MVP.
It had the lowest immediate resurrection and highest benign survival AUC.
This result motivated the stricter conditional protocol.

\begin{table}[h]
\centering
\caption{Benign local MVP main-track result. HLC appears strongest before immediate matching and answer-bearing stress.}
\label{tab:benign-main}
\begin{tabular}{lccc}
\toprule
Method & Immediate $r@0$ & AUC survival & $h50$ \\
\midrule
NPO & 0.833 & 0.167 & 0 observed \\
Strong NPO & 0.222 & 0.778 & $>32$ censored \\
NPO + paraphrases & 0.167 & 0.833 & $>32$ censored \\
Syntactic diversification & 0.444 & 0.556 & mixed \\
HLC-SG & 0.056 & 0.944 & $>32$ censored \\
\bottomrule
\end{tabular}
\end{table}


\subsection{Answer-Bearing Stress Changes the Conclusion}

Under immediate-matched answer-bearing stress, current HLC variants did not beat the strongest local baselines.
The held-out stress table is the clearest locked diagnostic result.

\begin{table}[h]
\centering
\caption{Held-out answer-bearing stress result from the locked local MVP. HLC beats NPO + paraphrases but not immediate-matched Strong NPO or syntactic diversification.}
\label{tab:heldout-stress}
\begin{tabular}{lccccc}
\toprule
Method & $r@0$ & $c_{32}$ & $c_{128}$ & $c_{512}$ & Cond. AUC$_{512}$ \\
\midrule
HLC-SG & 0.056 & 0.037 & 0.685 & 0.978 & 0.351 \\
NPO + paraphrases & 0.167 & 0.733 & 1.000 & 1.000 & 0.057 \\
Strong NPO matched & 0.056 & 0.000 & 0.263 & 0.978 & 0.410 \\
Syntactic diversification matched & 0.056 & 0.056 & 0.404 & 0.981 & 0.480 \\
\bottomrule
\end{tabular}
\end{table}


\subsection{GPT-2 LoRA Smoke Confirms the Larger-Model Path}

The GPT-2 LoRA lane validates the training and evaluation path but does not establish a method win.
K2 improves some short-horizon signals over K1, but still trails Strong NPO by $t=128$.

\begin{table}[h]
\centering
\caption{GPT-2 LoRA long-horizon diagnostic. K2 improves over K1 at $t=128$ but does not beat Strong NPO.}
\label{tab:gpt2-long}
\begin{tabular}{lccccc}
\toprule
Method & $r@0$ & $c_{32}$ & AUC$_{32}$ & $c_{128}$ & AUC$_{128}$ \\
\midrule
HLC-SG K1 GPT-2 & 0.000 & 0.278 & 0.920 & 0.833 & 0.529 \\
HLC-SG K2 GPT-2 & 0.000 & 0.278 & 0.920 & 0.611 & 0.584 \\
Strong NPO & 0.000 & 0.333 & 0.906 & 0.556 & 0.588 \\
\bottomrule
\end{tabular}
\end{table}


\subsection{Qwen3-0.6B Seeded Long-Horizon Gate}

The Qwen3-0.6B lane is the active decision gate for a tuned HLC K4 candidate against immediate-matched constrained gradient ascent.
The unseeded/provisional long-512 result should not be used for paper claims.
The seeded rerun must populate this table.

\begin{table}[h]
\centering
\caption{Qwen3-0.6B seeded long-horizon decision gate. Fill from the seeded Qwen report after the rerun completes.}
\label{tab:qwen-pending}
\small
\begin{tabular}{lccccc}
\toprule
Method & $r@0$ & $c_{32}$ & $c_{128}$ & $c_{512}$ & Cond. AUC$_{512}$ \\
\midrule
HLC K4 & \todo{fill} & \todo{fill} & \todo{fill} & \todo{fill} & \todo{fill} \\
Matched GA & \todo{fill} & \todo{fill} & \todo{fill} & \todo{fill} & \todo{fill} \\
\bottomrule
\end{tabular}
\end{table}


\paragraph{Current interpretation.}
If tuned HLC beats matched gradient ascent on conditional resurrection or conditional AUC after the seeded rerun, the paper can include a positive Qwen diagnostic result.
If it loses, the paper should remain a rigorous negative result centered on the evaluation protocol.


\section{Discussion}
The main lesson is that durable unlearning requires a different evidence standard than immediate forgetting.
In the local MVP, a method that looked best on benign survival curves failed once the comparison controlled for $r@0$ and used answer-bearing stress.
This is exactly the kind of failure the harness is meant to expose.

The results also suggest that the relearning pool matters.
Benign pools test stability under ordinary updates, but they are too weak to establish durable forgetting.
Answer-bearing pools are harsher and more informative because they probe whether a suppressed fact can be reintroduced through plausible post-unlearning updates.

The Qwen experiments add a second lesson.
When Strong NPO could not match HLC's immediate forgetting, constrained gradient ascent provided a useful frontier baseline.
This baseline is not SOTA in the broader literature, but it is locally decisive: if HLC cannot beat an immediate-matched constrained GA baseline, it cannot support a strong method claim in this harness.

The strongest paper angle is therefore methodological.
The harness gives reviewers concrete evidence that the authors are not mistaking aggressive forgetting for durability.
That is valuable even if the current HLC objective remains negative.


\section{Limitations and Next Evidence Gates}
This draft has several limitations.

First, the current strongest evidence is local and diagnostic.
The Qwen3-0.6B lane is modern enough to be useful, but it is not a full SOTA benchmark.
Larger models, more seeds, and standardized OpenUnlearning-style baselines would be needed for a broad claim.

Second, current stress pools are synthetic biography style.
They are controlled and reproducible, but future work should test more diverse domains and forgetting types.

Third, the HLC objective explored here may be too shallow.
K1/K2/K4 variants can improve short-horizon behavior, but prior locked results show that improvements do not always persist to $t=128$ or $t=512$.

Fourth, the current draft intentionally avoids external related-work citations.
The next writing pass should add exact citations for TOFU, NPO, machine unlearning, model editing, and durable/privacy-oriented unlearning benchmarks.

The next evidence gates are:

\begin{enumerate}[leftmargin=*]
  \item finish seeded Qwen3-0.6B long-512 HLC vs matched GA;
  \item add or confirm a third Qwen seed if compute allows;
  \item rerun final aggregate tables from canonical scripts;
  \item decide whether the results section is positive HLC evidence or a negative diagnostic study;
  \item add related-work citations and final figures.
\end{enumerate}


\section{Conclusion}
Durable unlearning should be evaluated as survival under relearning, not only as immediate forgetting.
This draft presents a local-first survival harness, immediate-matched baseline selection, and conditional resurrection metrics that make that distinction operational.
\method{} provides a concrete half-life-control objective for testing whether simulated relearning penalties can improve durability.

The current locked evidence does not justify a SOTA claim.
It does justify a rigorous paper narrative: naive survival curves can overstate durability, and immediate-matched answer-bearing stress is a stronger test.
The final Qwen seeded long-horizon result will decide whether the tuned HLC candidate becomes a positive method result or remains a useful negative case study.


\end{document}
